% correctness.tex — \input'd from experiments/experiments.tex
%
% Status: APPROVED (Jörn 2026-02-16)
%
% Sources: experiments/correctness/correctness.rs (dataset generator + tests),
%   experiments/correctness/correctness.jsonl (generated dataset),
%   experiments/correctness/README.md (experiment design documentation).
%
% This section documents the mathematical propositions verified to establish
% correctness of the capacity computation implementation.

\subsection{Correctness Verification}
\label{sec:correctness}

We test our implementation on several propositions that must hold for any
algorithm that computes the correct capacity.
Table~\ref{tab:correctness-algorithms} lists the algorithm variants we use.
Table~\ref{tab:verification-capacity-counts} shows the 6 tests we run, and
how many new capacity computations each test required. Tests reuse capacity
computations from other tests for development efficiency.

\begin{table}[h]
\centering
\begin{tabular}{ll}
\toprule
Algorithm & Description \\
\midrule
Pruned & HK2017 with adjacency-based orbit pruning (production) \\
Unpruned & HK2017 without pruning (validation only) \\
Billiard & Billiard algorithm (Lagrangian products only) \\
\bottomrule
\end{tabular}
\caption{Algorithm variants used for correctness verification.}
\label{tab:correctness-algorithms}
\end{table}

\begin{table}[h]
\centering
\small
\begin{tabular}{lrrr}
\toprule
Test & Pruned & Unpruned & Billiard \\
\midrule
Test 1: Direct comparison & 10 & 10 & 5 \\
Test 2: Literature & 7 & 0 & 4 \\
Test 3: Conformality & 10 & 0 & 5 \\
Test 4: Symplectic invariance & 10 & 0 & 0 \\
Test 5: Continuity & 10 & 0 & 0 \\
Test 6: Monotonicity & 0 & 0 & 0 \\
\midrule
\textbf{Total} & \textbf{47} & \textbf{10} & \textbf{14} \\
\bottomrule
\end{tabular}
\caption{Required capacity computations for verification. Test~1 generates
  10~base polytopes (5~random generic, 5~Lagrangian products). Tests~3--5
  reuse these base polytopes (scaled, transformed, perturbed). Test~6 uses
  existing capacities from tests~1--5. Total: 71~capacity values.}
\label{tab:verification-capacity-counts}
\end{table}

The 6 test propositions are:
\begin{enumerate}
  \item \textbf{Direct comparison} between all applicable algorithms on
    5~random polytopes and 5~random Lagrangian product polytopes.

  \item \textbf{Literature agreement} for all efficient (Pruned, Billiard)
    algorithms on the known published polytopes from
    Table~\ref{tab:literature-polytopes}.

  \item \textbf{Conformality}: $c(\alpha K) = \alpha^2 c(K)$ holds for the
    10~polytopes from test~1, with a random scaling~$\alpha$ each.

  \item \textbf{Symplectic invariance}: $c(MK) = c(K)$ holds for the
    10~polytopes from test~1, with a random~$M \in \mathrm{Sp}(4)$ each.

  \item \textbf{Continuity}: For the 10~polytopes from test~1, a small
    random perturbation of the facet normals produces a small change in
    capacity.

  \item \textbf{Monotonicity}: $\alpha^2 c(K_1) \le c(K_2)$ where we pick
    $\alpha$ maximal such that~$\alpha K_1 \subseteq K_2$, for all
    polytopes~$K_1, K_2$ in the dataset.
\end{enumerate}

\noindent
Volumes are computed via convex hull triangulation
(\texttt{qhull}), implemented in \texttt{geom::volume}.

\begin{table}[h]
\centering
\small
\begin{tabular}{llrrrr}
\toprule
Polytope & $F$ & $c_{\mathrm{EHZ}}$ & $\mathrm{vol}(K)$ & $\mathrm{sys}(K)$ & Source \\
\midrule
Simplex & 5 & 0.25 & 0.0417 & 0.750 & Nir 2013 \\
Hypercube $[-1,1]^4$ & 8 & 4.0 & 16.0 & 0.500 & HK2019 Ex~4.6 \\
HK-O pentagon & 10 & 3.441 & 5.657 & 1.047 & HK-O 2024 Prop~1.4 \\
Lag.\ $\triangle \times \triangle$ & 6 & 1.5 & 1.6875 & 0.667 & HK2017 + billiard \\
Sym.\ $\triangle \times \triangle$ & 6 & 1.299 & 1.6875 & 0.500 & Moser's theorem \\
Lag.\ $\triangle \times \square$ & 7 & 1.5 & 1.2990 & 0.866 & HK2017 + billiard \\
Sym.\ $\triangle \times \square$ & 7 & 1.0 & 1.2990 & 0.385 & Moser's theorem \\
\bottomrule
\end{tabular}
\caption{Known polytopes from the literature used in test~2.}
\label{tab:literature-polytopes}
\end{table}
